ReadProof-Bench is specified with four tracks (Figure~\ref{fig:bench}, Table~\ref{tab:tracks}). This version
evaluates the two that measure acquisition and channel coverage directly, Track~A and Track~K; Tracks~B and~C are
specified, and their tooling is released, with results in a later version. Every item records its collection
date and every result records the profile, the declared user agent and the network origin, because posts
disappear quickly and selectively \citep{pew2024disappears,zhu2013velocity} and re-collected social-media data
drift \citep{elmas2023persistence}.

\begin{table}[t]
\centering\small
\caption{The four tracks of ReadProof-Bench and what this version evaluates.}
\label{tab:tracks}
\begin{tabular}{@{}l p{2.7cm} p{4.0cm} p{3.0cm} p{1.9cm}@{}}
\toprule
Track & Unit & Gold & Headline metrics & This version \\
\midrule
A access & a link form of a public post & a reference fingerprint of the post (its search-result title) & coverage per
  form, answer-from-wall, honesty penalty & 145 posts, 522 forms \\
K canaries & one channel of one of the author's posts & the planted code, exact by construction & correct codes
  per channel, fabrication & 7 posts, 28 items \\
B evidence & a question about one post & answer, channel, locator, verbatim quote & accuracy per channel, quote
  validity & specified \\
C counting & a set of posts on one topic & per-post labels & count error, interval coverage & specified \\
\bottomrule
\end{tabular}
\end{table}

\paragraph{Track A: access across link forms.}
We collected 145 public posts on nine platforms (Douyin, Xiaohongshu and RedNote, Bilibili, YouTube, TikTok,
Weibo, Toutiao, X and WeChat articles) on 21 November 2026 with a web search engine, about fifteen per platform
spread over eight topic domains; no platform was crawled to build the list. Each seed keeps the URL exactly as
the engine returned it, the post identifier, the topic and language, and the title or snippet the engine showed,
which serves as the post's reference fingerprint. For each post the platform module emits up to four link forms
that people actually paste (Appendix~\ref{app:forms}): the canonical desktop address, the mobile address, the
share text in the app's own template with the address embedded, and the most informative platform variant (a share token stripped
for Xiaohongshu, the share path for Toutiao, the mobile detail page for Weibo, \texttt{youtu.be} for YouTube),
522 post--form items in all. A system \emph{serves} an item when at least 60\% of the fingerprint's character
bigrams (Chinese) or words (other languages) occur in the text it obtained. The scoring unit is the link form,
not the post, because the same post can open from one form and fail from another. Because the fingerprint is in the
post's language, we also count a page that holds the post in other words, such as a machine translation
(Section~\ref{sec:setup}). Weibo's indexed titles are
summaries rather than the post's words, so for Weibo we also report coverage against an audited capture of each
post's opening words.

\paragraph{Track K: planted canaries.}
The author published seven posts on 21 November 2026 on the author's own accounts: videos on Douyin, Bilibili,
YouTube, TikTok, RedNote and X, and a four-image note on RedNote. Each post carries distinct random codes, each
in exactly one channel: spoken only (two nouns and a number, masked in the burned-in subtitle), on screen only
(never spoken), in the caption only, in the author's first comment only (pinned where the platform's web
interface allows it), or inside one carousel image. The image note also has a spoken code that was never planted,
which makes ``what code is spoken in this post?'' an absent-channel control: the only correct answer is that
there is none. A correct code proves that the system read that channel on that platform, without access to its
logs; a different code is a fabrication; ``not found'' is an abstention. Codes differ across posts and channels,
so none can be guessed. The answer key was committed by its SHA-256 hash before any system ran and is published
after the evaluation rounds. Spoken Chinese codes are scored by toneless pinyin, so a homophone spelling
(for example \zh{海豚} heard as \zh{海屯}, both \emph{haitun}) still proves that the audio was transcribed. Track~K is the one slice whose media
can be released in full, because the author holds the rights. It extends the planted codes of the robots.txt
study of AI assistants \citep{lopezfonseca2026robots} to third-party platforms that cannot be instrumented from
outside.

\paragraph{Tracks B and C (specified).}
Track~B asks questions answerable from exactly one channel of one third-party post, with gold from the agreement
of two independent extractors, a closed-book filter and an answerability filter \citep{lee2026vragbench,
wei2025browsecomp}. Track~C asks ``of these $M$ posts, how many say $A$?'' over sets of posts on one topic,
with per-post labels and one verified quote per counted post, as QQSUM derives prevalence from matches on
individual comments \citep{tang2025qqsum}. Both are specified in the released design document and are left to
the next version.

\paragraph{Release.}
Third-party content is released only as post identifiers, URLs, titles and hashes; no page text and no media.
Track~K media and the answer key are released in full once the evaluation rounds end. The systems in this version read the links they are
given and use no web search, so the answer key cannot leak to them through search. Code is released under
Apache-2.0 and data under CC~BY~4.0; Appendix~\ref{app:datasheet} is the datasheet \citep{gebru2021datasheets}.
